% Full proof transcription of work/imprimitive/IMPRIMITIVE_PROOF.md.
% 2026-11-21; source author /root/imprimitive; transcription /root/g34.
\section{The standard imprimitive representations factor through cyclic Pryms}
\label{sec:HD-I}

Throughout this section $d\in\{3,4,6\}$, $e\mid d$, $m\ge9$, and
$K=\mathbb Q(\zeta_d)$. The representation under consideration is the
standard $m$-dimensional reflection representation of $G(d,e,m)$.
We prove that every associated signature-$(n,n)$ period map has rank at
most $3n$, with the bound $3n-1$ for $e<d$. For $n=5$ the maximum image
dimension over all branch components of a fixed group is exactly $15$
for $e=d$ and exactly $14$ for $e<d$. The exact generic ranks of individual
branch components are not classified here.

\paragraph{HD-I01 (the inducing morphism).}
Use the convention
\[
G(d,e,m)=\left\{\operatorname{diag}(\zeta_d^{a_1},\ldots,
\zeta_d^{a_m})P_\pi:\sum_i a_i\equiv0\pmod e\right\}.
\]
Let $V=K^m$, let $H$ stabilize the line $Ke_1$, let
$\chi:H\to\mu_d$ be the scalar action on that line, and put
$N=\ker\chi$. Then
\[
K[G]\otimes_{K[H]}K_\chi\longrightarrow V,
\qquad g\otimes1\longmapsto ge_1
\]
is an isomorphism of $G$-representations. The representation $V$ is
absolutely irreducible, has character field $K$, and is different from
its complex conjugate.

\paragraph{Proof.}
The element $\operatorname{diag}(\zeta_d,\zeta_d^{-1},1,\ldots)$
belongs to $H$ and proves that $\chi$ is onto. The displayed map respects
the balanced tensor relation. The $G$-orbit of $Ke_1$ consists of the
$m$ coordinate lines, so $[G:H]=m$. The map is onto, and its source and
target both have dimension $m$.
The coordinate characters of the diagonal subgroup are distinct: a
third coordinate permits an exponent $1$ at either of a specified pair
and exponent $-1$ at that third coordinate. Thus an invariant subspace
is a sum of coordinate lines. The contained symmetric group acts
transitively on those lines, proving irreducibility after any extension
of $K$. The diagonal element with exponents $(1,1,-2,0,\ldots)$ has
nonreal trace for each $d=3,4,6$. It belongs to every group in question.
The character values lie in $K$, and a nonreal element generates this
quadratic field. Hence $V$ and $\bar V$ differ and their rational simple
summand is $\operatorname{Res}_{K/\mathbb Q}V$.
For any $K[G]$-module $U$, the corresponding Frobenius map is explicitly
\[
\operatorname{Hom}_{K[G]}(V,U)
\longrightarrow\{u\in U:hu=\chi(h)u\text{ for every }h\in H\},
\qquad F\longmapsto F(e_1).
\]
Its inverse sends $u$ to the map $ge_1\mapsto gu$; well-definedness is
exactly the induction isomorphism just proved.

\paragraph{HD-I02 (quotients and the polarized isogeny).}
Let $Z\to\mathbb P^1$ be a connected smooth $G$-Galois cover and put
\[
X=Z/N,\qquad Y=Z/H,\qquad q:Z\to X,\qquad f:X\to Y.
\]
The primitive cyclic $d$-cover Prym of $f$ is isogenous to the
multiplicity factor for $V$ in $J(Z)$, through a specified polarized
isogeny. In families this gives a factorization of their period maps
through the moduli of cyclic covers $f$.

\paragraph{Proof.}
The cover $f$ is connected and cyclic of degree $d$, while
$Y\to\mathbb P^1$ has degree $m$. Use the cohomological left action
$g\mapsto(g^{-1})^*$. Pullback identifies
$H^1(X,\mathbb Q)$ with $H^1(Z,\mathbb Q)^N$; the averaging
correspondence is $|N|^{-1}q^*q_*$. Composing with HD-I01 identifies
$\operatorname{Hom}_G(V,H^1(Z,K))$ with the $\chi$-eigenspace of
$H^1(X,K)$, and identifies their inverse-character spaces as well.
Pullback, norm and sums of automorphism graphs preserve the Hodge
filtration.
More precisely, put
\[
p=\frac1{|H|}\sum_{h\in H}
\bigl(\chi(h^{-1})+\bar\chi(h^{-1})\bigr)h\in\mathbb Q[H].
\]
The coefficients are rational because $\zeta_d^a+\zeta_d^{-a}$ is an
integer for the three values of $d$. The two character projectors are
orthogonal idempotents. Frobenius reciprocity in HD-I01 shows that $p$
vanishes on every rational simple constituent except
$\operatorname{Res}_{K/\mathbb Q}V$ and has rank one over $K$ on that
constituent. It is self-adjoint under $h\mapsto h^{-1}$.
Let $P\subset J(X)$ be the primitive-character abelian subvariety and
let $B=\operatorname{im}(|H|p:J(Z)\to J(Z))$. The identity
$q^*e_{\rm primitive}=pq^*$ gives $q^*P=B$. Consequently
\[
\alpha=q^*|_P:P\longrightarrow B
\]
is an isogeny, and norm/pullback duality under the canonical Jacobian
polarizations gives
\[
\alpha^*(\Theta_Z|_B)=|N|(\Theta_X|_P).
\]
Indeed $q_*q^*=[|N|]$, which also shows
$\ker\alpha\subseteq P[|N|]$.
Another fixed primitive matrix idempotent in the same component gives
a fixed group-algebra isogeny. Its restricted polarization changes
by the positive rational value of the standard invariant Hermitian
form on the corresponding line; the displayed self-adjoint $p$ suffices
throughout.
In a Hurwitz family these correspondences are relative algebraic
correspondences. After finite level and restriction to a connected
component, their finite kernels are locally constant. Marked periods
are related by a fixed rational symplectic similitude. Thus the period
map factors through the cyclic-cover moduli followed by a fixed
isogeny correspondence, whose differential is invertible.

\paragraph{HD-I03 (all branch signatures as colored-cycle conditions).}
For a branch monodromy $\gamma_b$, let $c$ run over the cycles of its
permutation part. Write $\ell_c$ for a cycle length and
$s_c\in\{0,\ldots,d-1\}$ for its exponent sum modulo $d$. Define
\[
L=\sum_{b,c}(\ell_c-1),\quad
r_s=\#\{(b,c):s_c=s\}\ (1\le s<d),\quad
r=\sum_s r_s,\quad S=\frac1d\sum_s sr_s.
\]
Then
\[
g(Y)=1-m+\frac L2,
\qquad
\{h^{1,0}_\chi,h^{1,0}_{\bar\chi}\}
=\{g(Y)-1+S,\ g(Y)-1+r-S\}.
\]
Signature $(n,n)$ is equivalent to
\[
r=2(n-g(Y)+1),\qquad
\sum_{s=1}^{d-1}(2s-d)r_s=0.
\]
The relevant Hurwitz families are exactly the braid orbits of generating
product-one tuples in $G(d,e,m)$ satisfying these equations.

\paragraph{Proof.}
The characteristic polynomial on a colored cycle is
$T^{\ell_c}-\zeta_d^{s_c}$, and its eigenangles are
\[
\frac{k+s_c/d}{\ell_c},\qquad 0\le k<\ell_c.
\]
Their sum is $(\ell_c-1)/2+s_c/d$; the sum of their negative fractional
parts is $(\ell_c-1)/2+\mathbf1_{s_c\ne0}-s_c/d$. A zero angle
occurs precisely when $s_c=0$.
Riemann--Hurwitz for $Y\to\mathbb P^1$ gives its stated genus. A cycle
$c$ gives one point of $Y$ over the branch value. The monodromy of
$X\to Y$ there is $\zeta_d^{s_c}$, obtained from
$\gamma_b^{\ell_c}$ acting on the selected coordinate. Thus $s_c\ne0$
records exactly a cyclic branch point of exponent $s_c$.
The cyclic eigensheaves proved in HD-I05, together with Riemann--Roch,
give the displayed multiplicities, equivalently the Chevalley--Weil
angle sums above. Equating them to $n$ gives the two equations.
In particular, the balance equations are
\[
\begin{array}{c|l}
d&\text{balance}\\\hline
3&r_1=r_2\\
4&r_1=r_3\\
6&2r_1+r_2=r_4+2r_5.
\end{array}
\]
For $d=4$ the count $r_2$ is unrestricted by balance, as is $r_3$ for
$d=6$. Each tuple member also obeys its exponent-sum condition modulo
$e$. Product one and generation remain part of the stated
characterization. Riemann existence realizes every such tuple; for
$b$ ordered branch values its source dimension modulo
$\operatorname{PGL}_2$ is $b-3$.

\paragraph{HD-I04 (universal rank bounds).}
Every signature-$(n,n)$ source above has period differential rank at
most $3n$. For $e<d$ its rank is at most $3n-1$. In particular the
standard imprimitive construction cannot dominate a full Weil domain
for $n\ge4$.

\paragraph{Proof.}
A cyclic cover of a genus-$g$ base with $r$ branch points and specified
exponents is determined, up to finitely many choices, by those pointed
curve data and a $d$th root of the weighted branch divisor. Hence its
moduli dimension is
\[
3g-3+r=g+2n-1\le3n,
\]
because $r=2(n-g+1)\ge0$. For $g=0,1$ the same formula applies:
$n\ge4$ supplies enough markings for stability. HD-I02 transfers this
bound to the original period map.
For $e<d$, the case $r=0$ would mean that every colored cycle of every
local monodromy has sum zero modulo $d$. Every generator would then
lie in the proper subgroup $G(d,d,m)$, contradicting full monodromy.
Therefore $r\ge2$, $g\le n$, and the bound improves to $3n-1$.
The target dimension is $n^2$ and $3n<n^2$ for $n\ge4$. For $n=5$
the two bounds are $15$ and $14$, against target dimension $25$.

\paragraph{HD-I05 (the primitive cyclic codifferential).}
Let $X\to Y$ be a smooth cyclic cover with branch exponents
$a_j\in\{1,\ldots,d-1\}$, reduced branch divisor $R$, and
$\mathcal L^d=\mathcal O_Y(\sum_j a_jP_j)$. Its two primitive
holomorphic eigenspaces are
\[
H^0(\omega_Y\mathcal L),\qquad
H^0(\omega_Y\mathcal L^{-1}(R)).
\]
The codifferential of its primitive Prym period map is, up to a
nonzero scalar, multiplication
\[
\mu:H^0(\omega_Y\mathcal L)\otimes
H^0(\omega_Y\mathcal L^{-1}(R))
\longrightarrow H^0(\omega_Y^2(R)).
\]

\paragraph{Proof.}
Normalization of $y^d=s$ gives the $i$th eigenline in $f_*\mathcal O_X$
as
$\mathcal L^{-i}(\sum_j\lfloor ia_j/d\rfloor P_j)$.
Finite-map duality gives the canonical eigensheaves. For $i=1$ all
floors vanish; for $i=d-1$ they are $a_j-1$, giving precisely the
two spaces displayed.
For the general-exponent differential formula, use Serre duality on
infinitesimal Torelli. In the local tame chart $x=t^k$, invariant
vector fields upstairs descend to $xh(x)\partial_x$. Invariant
deformations are thus $H^1(T_Y(-R))$. Dualizing the cup action on
inverse eigenspaces gives multiplication followed by trace. Its product
has at most a simple pole at each branch point, hence is the stated
section of $\omega_Y^2(R)$, and trace scales by the nonzero degree
$d$. This proves the required general-exponent version.
Lange--Ortega, \emph{Prym varieties of cyclic coverings}, Geometriae
Dedicata 150 (2011), 391--403, Proposition 4.1 and Lemma 4.2,
give the multiplication formulation under their reduced-branch
convention; the author preprint arXiv:0805.1020 has the proof on
printed pp.~7--10.

\paragraph{HD-I06 (rank 15 from an unramified genus-six source).}
For each $d=3,4,6$ there is a smooth unramified connected cyclic
$d$-cover of a genus-six curve for which the primitive period
differential has rank exactly $15$.

\paragraph{Proof.}
Begin with $Y_0$ obtained from $\mathbb P^1$ by identifying the six
pairs $a_i=2i$, $b_i=2i+1$, $0\le i\le5$. This is an irreducible
stable genus-six curve. Let $\eta$ be trivial on its normalization,
with gluing $\zeta_d$ at each node. Its order is exactly $d$: an
isomorphism trivializing a power is a constant automorphism on the
normalization and requires $\zeta_d^k=1$. The associated cyclic
torsor is connected because a node loop maps onto $\mu_d$.
Set
\[
D(z)=\prod_{j=0}^{11}(z-j),
\]
\[
F_i(u)=\frac{D(z)}{z-a_i}-u\frac{D(z)}{z-b_i}
       -\frac{D(z)}{z-a_5}+u\frac{D(z)}{z-b_5},
\qquad 0\le i\le4.
\]
The forms $F_i(\zeta_d)\,dz/D$ and
$F_i(\zeta_d^{-1})\,dz/D$ are bases of
$H^0(\omega\eta)$ and $H^0(\omega\eta^{-1})$.
Indeed their node residues obey $r_{b_i}=-ur_{a_i}$; regularity at
infinity gives $(1-u)\sum_i r_{a_i}=0$. The five residue differences
are a basis of this kernel.
Their products have denominator $D^2$, numerator degree at most $20$,
and equal double-pole coefficients at each pair, so belong to
$H^0(\omega_{Y_0}^2)$. This space has dimension $15$: the normalization
divisor has degree $20$ and six matching conditions, or nodal
Riemann--Roch and $H^1(\omega^2)=0$ give $3g-3$.
Choose the fifteen products $F_i(\zeta_d)F_j(\zeta_d^{-1})$,
$0\le i\le j\le4$, in lexicographic order, and take coefficients of
$z^0,\ldots,z^{14}$. The exact finite-field certificate gives
\[
\begin{array}{c|c|c}
d&\zeta_d\bmod61&\det\bmod61\\\hline
3&13&34\\4&11&18\\6&14&51.
\end{array}
\]
All determinants are nonzero, proving independence over $K$ and rank
$15$. The exact verification script reconstructs the polynomials and
stores the matrices in its JSON output; its filename is
\texttt{verify\_\allowbreak period\_\allowbreak ranks.py}.

The stable curve has a smoothing: Stacks Project, Theorem 109.25.3,
tag 0E9C, states smoothness of the stable-curve stack and density of its
smooth locus. The line bundle $\eta$ with its $d$th-power trivialization
extends locally over this smoothing. Line-bundle obstructions lie in
$H^2(\mathcal O_{Y_0})=0$, and multiplication by $d$ is invertible on
the infinitesimal discrepancy in the power isomorphism. Equivalently,
its finite etale $\mu_d$-torsor lifts locally. The torsor is locally a
disjoint union at each node; smoothing therefore introduces no node
inertia. Exact order and connectedness persist after shrinking.
The dimensions of $H^0(\omega\eta^{\pm1})$ remain $5$, because
$\eta^{\pm1}$ remain nontrivial of degree zero; the target dimension
remains $15$. Cohomology and base change makes multiplication a map
of vector bundles, and the nonzero minor persists on smooth nearby
fibres. HD-I05 gives the claimed smooth period rank.

\paragraph{HD-I07 (rank 14 with two branch points).}
For every proper divisor $e$ of $d=3,4,6$ there is a connected cyclic
$d$-cover of a smooth genus-five curve with two branch exponents
$e,d-e$ for which the primitive period differential has rank exactly
$14$.

\paragraph{Proof.}
On $\mathbb P^1$, mark $P=0$, $Q=\infty$, and identify five pairs
\[
a_i=7^{d(2i+1)},\qquad b_i=7^{d(2i+2)},\qquad 0\le i\le4.
\]
Let $\mathcal L$ normalize to $\mathcal O(\infty)$ with gluing
$\lambda=\zeta_d7^e$ at each node. The section $s=z^e$ of
$\mathcal L^d$ descends because $s(b_i)=\lambda^d s(a_i)$.
Its divisor is $eP+(d-e)Q$. On the normalization the cyclic cover
has $e$ components when $e>1$, and gluing advances the root label by
one: $b_i^{e/d}=7^e a_i^{e/d}$ and the extra factor is $\zeta_d$.
All components are joined by these nodes. Thus the cover is connected
and unramified at the nodes.
Put
\[
D=\prod_i(z-a_i)(z-b_i),\qquad Q_c=D/(z-c),
\]
\[
A_i=Q_{a_i}-\lambda Q_{b_i},\qquad
B_i=z(Q_{a_i}-\lambda^{-1}Q_{b_i})-(1-\lambda^{-1})D.
\]
The forms $A_i\,dz/D$ are a basis of $H^0(\omega\mathcal L)$:
a simple pole at infinity is allowed, so the five node residues are
free. The forms $B_i\,dz/(zD)$ are a basis of
$H^0(\omega\mathcal L^{-1}(P+Q))$; the added residue at zero enforces
regularity at infinity, and the inverse gluing relation holds at each
node. Products have numerators $A_iB_j$ and common denominator $zD^2$.
They belong to $H^0(\omega^2(P+Q))$, whose dimension is $14$:
the normalization divisor has degree $18$, giving nineteen
coefficients subject to five matching conditions.
Select the following fourteen ordered pairs:
\[
\begin{split}
&(0,0),(0,1),(0,2),(0,3),(0,4),\\
&(1,0),(1,1),(1,2),(1,3),(1,4),\\
&(2,0),(2,2),(3,3),(4,4).
\end{split}
\]
Take coefficients of $z^0,\ldots,z^{13}$. Modulo $241$, where $7$ is
primitive and $\zeta_d\mapsto7^{240/d}$, the determinants are
\[
\begin{array}{cc|rrr}
d&e&\zeta_d&\lambda&\det\\\hline
3&1&15&105&150\\
4&1&177&34&152\\
4&2&177&238&15\\
6&1&16&112&43\\
6&2&16&61&166\\
6&3&16&186&213.
\end{array}
\]
All entries belong to the reduction of
$\mathbb Z[1/7,\zeta_d]$, so the nonzero minors prove
characteristic-zero rank $14$.
The marked curve and the root $(\mathcal L,s)$ smooth: line-bundle
deformations are unobstructed, and multiplication by $d$ corrects
the root isomorphism as in HD-I06. The marked sections remain disjoint
and smooth. Normalization of the cyclic cover is smooth over smooth
base fibres; the central node-unramified cover is connected, and
connectedness persists. Cohomology and base change preserves the
source dimensions $5,5$ and target dimension $14$. Thus the nonzero
minor persists to smooth primitive period maps.

\paragraph{HD-I08 (full monodromy and the exact maxima).}
For $n=5$ and every fixed $d,e,m$ in the stated range, the maximum
period-image dimension over branch components with full monodromy
$G(d,e,m)$ equals $15$ for $e=d$ and $14$ for $e<d$.

\paragraph{Proof.}
A smooth curve $Y$ of genus $g\le6$ has a very ample line bundle of
every degree $m\ge g+3$. To prove the needed general-line-bundle form,
choose a general nonspecial line bundle $\mathcal M$ of degree $m$.
Failure to separate a length-two divisor $D_2$ forces
$H^1(\mathcal M(-D_2))\ne0$, so
$K_Y\mathcal M^{-1}(D_2)$ is effective of degree $2g-m$.
The failure locus in $\operatorname{Pic}^mY$ is therefore contained
in the image of
$\operatorname{Sym}^{2g-m}Y\times\operatorname{Sym}^2Y$, of dimension
at most $2g-m+2\le g-1$. A negative symmetric power gives the empty
locus. The nonspeciality failure locus likewise has smaller dimension.
Thus a general $\mathcal M$ is very ample.
A general pencil in this embedding is basepoint-free and gives a
simply branched degree-$m$ map $h:Y\to\mathbb P^1$ with distinct branch
values. The exact published input is SGA 7 II, Expose XVII,
Definition 2.2 and Theorem 2.5.2: every smooth proper embedding in
characteristic zero is a Lefschetz embedding. For a curve its pencil
axis is disjoint, and every singular hyperplane section has precisely
one ordinary double intersection. The finitely many marked points
can also be required to have distinct images away from these branch
values, by the corresponding open conditions.
Apply this with $g=6$ in HD-I06 and $g=5$ in HD-I07; $m\ge9$ suffices.
The monodromy of $X\to Y\to\mathbb P^1$ acts on $m$ blocks of $d$
sheets and lies in $G(d,1,m)$. A simple branch of $h$ gives a colored
transposition
\[
s_{ij}(a):(i,k)\mapsto(j,k+a),\qquad(j,k)\mapsto(i,k-a).
\]
The two cyclic branch points, when present, give diagonal rotations
by $e$ and $-e$ on individual blocks. In the unramified case put $e=d$.

The transposition graph is connected because the monodromy of $h$
is transitive. Choose a spanning tree. Diagonal conjugation kills its
edge labels, so the group contains the ordinary $S_m$ generated by
the tree transpositions. A further colored transposition times its
ordinary transposition gives $a(E_i-E_j)$ in the diagonal subgroup.
Conjugating by $S_m$ gives all such differences. The branch rotation
gives $eE_i$ for every $i$. Put
\[
q_0=\gcd(d,e,\text{all remaining edge labels}).
\]
The generated diagonal subgroup is exactly
\[
\{v\in(\mathbb Z/d)^m:
 v_i\equiv0\pmod{q_0}\text{ for all }i,\quad
 \sum_i v_i\equiv0\pmod e\}.
\]
Every generator satisfies these restrictions. Conversely Bezout
combinations give $q_0(E_i-E_j)$; together with $eE_i$ they produce
every displayed vector, by moving its coordinates to the first
coordinate and then using the allowed total sum. There are exactly
$q_0$ orbits on the $md$ sheets: the sheet label modulo $q_0$ is
preserved, and the difference vectors move a chosen label by $q_0$.
The composite cover is connected, so its monodromy is transitive.
Therefore $q_0=1$, and the full monodromy is $G(d,e,m)$.

The auxiliary pencils form an open family over the cyclic-cover
moduli. Forgetting the pencil is dominant and has surjective
differential at a general point in characteristic zero. This also
follows from the relative Picard and Grassmann parameter spaces of
the nonspecial line bundles and their pencils. A finite marking of
the Galois closure preserves rank. HD-I02 identifies the period
differential with the cyclic Prym differential after this forgetful
map and an invertible isogeny differential. The ranks $15$ and $14$
proved in HD-I06 and HD-I07 are attained on Hurwitz components with
full group $G(d,e,m)$. HD-I04 supplies the matching upper bounds.

\paragraph{Framework and audit.}
The quotient maps and the group-algebra idempotents give algebraic
correspondences on degree-one relative motives, and the pullback/norm
maps give the same maps on Betti and $\ell$-adic realizations. They
spread over a model with $d|G|$ and the needed discriminants inverted.
At the realization level exterior powers are determinant
$\lambda$-operations. On the $H^1$ tensor piece of a curve power,
geometric permutation includes its Koszul sign: the geometric
symmetrizer realizes the usual exterior projector. This fixes the
sign convention for the corresponding motivic projector.
The smooth proper $\ell$-adic summands have weight one. The exact
polarized cyclic factorization is the source-dimension obstruction
in the supplied ATTACK Proposition O and Corollary O$''$ framework.
The programme's detecting-cycle condition remains condition A,
cote 16 scan 050, as corrected in SIGN\_NEXT SN-C03/C04; the
projector belongs to the correspondence framework of scans
048--049. These statements produce the indicated period-image bounds.
They leave the full $n^2$-dimensional detecting-family gate open for
this imprimitive route.

The author read the whole proof and reran the standalone certificate
with its separate SymPy integer-polynomial and determinant check.
An independent agent read HD-I01--HD-I08 in full and found no fault
in the quotient, polarization, gluing, deformation, codifferential,
or full-monodromy arguments. Its separate rational-function evaluation
calculation at primes $109$ and $1009$ gives ranks $15$ and $14$
for all nine cases, using no imported coefficient matrices.
The precise events and both ordinary exact certificates are retained
in the accompanying files.
